% !TEX root = ../paper.tex

\section{Alderman Fixed Effects and Index Validation}
\label{sec:appendix_alderman_fe}

This appendix provides additional detail on how the aldermanic stringency index is constructed and validated.

\subsection{Construction of the Stringency Index}
\label{subsec:stringency_index_construction}

Aldermanic stringency measures whether comparable high-discretion permits take systematically longer under a particular alderman.
It uses the permit sample described in Section~\ref{subsec:permits_sample} through December 2022.
Permits must take at least one day to enter the processing-time comparison, although same-day permits still count toward the ward's recent permit volume.

\paragraph{Stage 1: Adjusting Processing Times.}
I first adjust processing times for observable reasons that otherwise could make permits slower in some wards.
The regression controls for the permit's distance to downtown and Lake Michigan, nearby CTA stations, ward demographics, recent ward permit volume, calendar month, permit type, and review type.
Appendix~\ref{app:data_sources} describes how block group demographics are allocated across ward boundaries to construct the ward controls.

\paragraph{Stage 2: Alderman Differences.}
I average the adjusted processing times by ward and month, then compare these averages across the aldermen serving each ward-month.
Formally, I regress the ward-month averages on alderman indicators $\alpha_{a(w,m)}$ and the same recent permit-volume control.
Ward-months are weighted by their permit counts.
Here $a(w,m)$ denotes the alderman serving ward $w$ in month $m$.
The resulting fixed effect is the persistent alderman-specific difference in adjusted processing time.

\paragraph{Empirical Bayes Shrinkage.}
Some aldermen have fewer permits and therefore less precise estimates.
I use empirical Bayes shrinkage to pull these noisier estimates toward the overall mean while leaving more precise estimates closer to their raw values:
$\hat{\alpha}_a^{\text{EB}} = B_a \hat{\alpha}_a$, where $B_a = \hat{\omega}^2/(\hat{\omega}^2 + \hat{\sigma}_a^2)$.
Here $\hat{\omega}^2$ is the estimated variance across aldermen and $\hat{\sigma}_a^2$ is the sampling variance of alderman $a$'s estimate.

\paragraph{Index Standardization.}
The stringency index, $S_a$, is standardized to have mean zero and unit standard deviation across aldermen.
Higher values indicate more stringent aldermen.

Table~\ref{tab:stage1_residualization} reports the adjustment regression used in the first stage.
The dependent variable is permit-level log processing time.
These controls remove predictable differences in timing across locations and permit types before I compare aldermen.
Panel B of Figure~\ref{fig:stringency_diagnostics} shows the distribution of the stringency scores.

\subsection{Validation of the Stringency Index}
\label{app:stringency_validation}

The stringency index should measure differences across aldermen rather than permanent differences across wards.
Panel C of Figure~\ref{fig:stringency_diagnostics} compares predecessor and successor aldermen within the same ward.
If fixed ward characteristics drove the scores, aldermen serving the same ward at different times would receive similar values.
The relationship is weak, which weighs against that explanation.

\begin{table}[htbp]
    \caption{Factors Associated with Permit Processing Times}
    \label{tab:stage1_residualization}
    \input{../tasks/create_alderman_uncertainty_index/output/stage1_regression_ptfeTRUE_rtfeTRUE_porchTRUE_cafeFALSE_2stage_volLAG1_BOTH_through2022.tex}

    \footnotesize
    \textit{Notes:} Permit-level linear regression used to construct the stringency index, using permit records through December 2022.
    Covariates are median household income, racial shares, homeownership, population, distance to downtown, distance to the lake, nearby CTA stations, and the previous month's permit volume.
    Fixed effects are listed in the table.
    Regressions are unweighted.
    Standard errors are clustered by ward and reported in parentheses. * $p<0.10$, ** $p<0.05$, *** $p<0.01$.
\end{table}

\begin{figure}[p]
    \centering
    \textbf{Panel A: Stringency maps in 2014 and 2022}\par
    \includegraphics[width=0.72\textwidth]{../tasks/strictness_score_map/output/uncertainty_score_map_comparison.pdf}

    \begin{minipage}[t]{0.43\textwidth}
        \centering
        \textbf{Panel B: Score distribution}\par
        \includegraphics[width=\linewidth]{../tasks/strictness_score_map/output/uncertainty_score_distribution_ptfeTRUE_rtfeTRUE_porchTRUE_cafeFALSE_2stage_volLAG1_BOTH_through2022.pdf}
    \end{minipage}\hfill
    \begin{minipage}[t]{0.43\textwidth}
        \centering
        \textbf{Panel C: Predecessors and successors}\par
        \includegraphics[width=\linewidth]{../tasks/strictness_score_map/output/uncertainty_score_predecessor_successor_ptfeTRUE_rtfeTRUE_porchTRUE_cafeFALSE_2stage_volLAG1_BOTH_through2022.pdf}
    \end{minipage}
    \caption{Variation in Aldermanic Stringency}
    \figurenotes{Stringency scores use permit processing times through December 2022 after adjusting for permit, location, ward, and calendar characteristics.
    Noisier alderman estimates are pulled toward the overall mean, and the resulting scores are standardized to mean zero and unit standard deviation.
    Panel A maps the scores of aldermen serving in January 2014 and January 2022 onto the corresponding ward boundaries; both maps use the same color scale, with warmer shades indicating more stringent aldermen.
    Panel B's dashed line marks the mean score of zero.
    Each point in Panel C pairs successive aldermen from the same ward.
    The dashed line indicates equal predecessor and successor scores; the solid line is a linear fit, with a shaded 95\% confidence interval for the fitted mean.}
    \label{fig:stringency_diagnostics}
    \label{fig:alderman_stringency_map}
    \label{fig:aldermanic_stringency_scores}
    \label{fig:predecessor_successor_stringency}
\end{figure}
